% 第5章 5.3.1：高斯核（sigma=0.15）RKHS 中目标 f=K_{0.2}-0.8K_{0.5}+0.9K_{0.8} 在观测点 Z={0.1,0.35,0.6,0.9} 上的投影。
% 数据 data.dat 由同目录 generate.py 计算：系数 a=(0.7729,0.0255,-0.3523,0.7455)，||f||^2=2.0392，||Pf||^2=1.2101。
% (a) f 与 Pf；(b) 余项 r=f-Pf，在观测点处为零。
\begin{tikzpicture}[font=\figurefont\footnotesize]
\pgfplotsset{every axis/.style={width=54mm,height=40mm,axis lines=left,xmin=0,xmax=1,
  xtick={0,0.1,0.35,0.6,0.9,1},xticklabels={$0$,$x_1$,$x_2$,$x_3$,$x_4$,$1$},
  x axis line style={draw=FigureStroke,line width=.6pt},y axis line style={draw=FigureStroke,line width=.6pt},
  tick style={draw=FigureStroke},tick label style={font=\figurefont\scriptsize,text=FigureInk},
  label style={font=\figurefont\scriptsize},xlabel={$t$},xlabel style={at={(axis description cs:1,0)},anchor=north west,yshift=2mm},
  xmajorgrids,ymajorgrids,grid style={draw=FigureGrid,line width=.4pt},
  title style={font=\figurefont\footnotesize},
  legend style={draw=none,fill=white,fill opacity=.85,text opacity=1,font=\figurefont\scriptsize,cells={anchor=west}}}}
\begin{axis}[name=A,ymin=-0.7,ymax=1.5,ytick={-0.5,0,0.5,1},title={(a) 目标与投影},legend style={at={(0.5,1.0)},anchor=north,legend columns=2}]
\addplot[draw=FigureBlue,line width=1.1pt] table[x=t,y=f] {figures/02-abstract-mathematics/03-function-spaces/tikz/c5-observation-residual/data.dat};
\addplot[draw=FigureGreen,line width=1.1pt,dashed] table[x=t,y=Pf] {figures/02-abstract-mathematics/03-function-spaces/tikz/c5-observation-residual/data.dat};
\addplot[only marks,mark=*,mark size=1.5pt,FigureInk,forget plot] coordinates {(0.1,0.7779) (0.35,0.1313) (0.6,-0.2420) (0.9,0.6978)};
\legend{$f$,$P_{M_Z}f$}
\end{axis}
\begin{axis}[at={(A.east)},anchor=west,xshift=9mm,ymin=-0.45,ymax=0.45,ytick={-0.4,-0.2,0,0.2,0.4},yticklabel style={/pgf/number format/fixed},title={(b) 余项$r=f-P_{M_Z}f$}]
\addplot[draw=FigureYellow,line width=1.1pt] table[x=t,y=r] {figures/02-abstract-mathematics/03-function-spaces/tikz/c5-observation-residual/data.dat};
\addplot[only marks,mark=*,mark size=1.5pt,FigureInk] coordinates {(0.1,0) (0.35,0) (0.6,0) (0.9,0)};
\end{axis}
\end{tikzpicture}
